% q_convention_bridge_appendix.tex
% q-convention bridge as cross-volume appendix.
%
% This appendix records the bridge identity q_KL^2 = q_DK with its
% normalization consequences. It is the canonical Vol I reference; Vol II appendix
% appendix_q_conventions_v2.tex and Vol III conventions.tex sec:q-convention-bridge
% point here.
%
% Standards: r-matrix discipline; hbar bridge identity; B-cycle monodromy;
% label uniqueness across three volumes; use degree or valence, not arity.

% Macro safety nets for cross-volume portability.
\providecommand{\bbC}{\mathbb{C}}
\providecommand{\bbZ}{\mathbb{Z}}
\providecommand{\bbQ}{\mathbb{Q}}
\providecommand{\bbR}{\mathbb{R}}
\providecommand{\fg}{\mathfrak{g}}
\providecommand{\hv}{h^{\vee}}
\providecommand{\Conf}{\mathrm{Conf}}
\providecommand{\KL}{\mathrm{KL}}
\providecommand{\DK}{\mathrm{DK}}
\providecommand{\GRT}{\mathrm{GRT}}
\providecommand{\Rep}{\mathrm{Rep}}

\chapter{The q-Convention Bridge}
\label{app:q-convention-bridge}

The multiplicative quantum parameter that controls braiding,
monodromy, and modular data of an affine Kac--Moody algebra at
non-critical level admits two standard normalisations. The
Kazhdan--Lusztig (KL) convention writes
$q_{\KL} := \exp(\pi i / (k + \hv))$; the Drinfeld--Kohno (DK)
convention writes $q_{\DK} := \exp(2\pi i / (k + \hv))$. They are
two normalisations of the same local monodromy. They are related by the squaring
identity
\begin{equation}
\label{eq:q-bridge}
\boxed{\; q_{\KL}^{2} \;=\; q_{\DK}. \;}
\end{equation}
The identity is structural, not cosmetic: it expresses the index-2
inclusion of the pure-braid subgroup into the full Artin braid
group on two strands, realised algebraically as a square root of
quantum-group parameters. The KL convention works on the universal
cover of the unordered configuration space (a framed datum); the
DK convention works on the base (an unframed datum); the square
root is the lift.

Theorem~\ref{thm:q-convention-bridge-main} fixes the relation through six
identities and three structural interpretations.  In particular, it prevents
the substitution \(\hbar=\pi i/(k+\hv)\) into
\(q=\exp(2\pi i\hbar)\), which would replace a root of unity by the real
decay factor \(\exp(-2\pi^2/(k+\hv))\).

\section{The bridge theorem}
\label{sec:q-bridge-theorem}

Throughout this appendix $\fg$ is a finite-dimensional simple Lie
algebra over $\bbC$ with normalised Killing form (long-root
squared length~$2$) and dual Coxeter number $\hv$. We work at
non-critical level $k \neq -\hv$, set
$\hbar := 1/(k + \hv)$ for the KZ coupling, and write $\Omega \in
\fg \otimes \fg$ for the Casimir tensor in the normalised Killing
form. The trace-form Casimir $\Omega_{\mathrm{tr}}$ is related to
$\Omega$ by
$\Omega = 2\hv \, \Omega_{\mathrm{tr}}$
(the standard normalisation fixed by long-root squared length~$2$).

\begin{theorem}[Q-convention bridge]
\label{thm:q-convention-bridge-main}
\ClaimStatusProvedHere
With notation as above, define
\[
q_{\KL} \;:=\; \exp(\pi i \hbar)
\;\;\in\;\; \bbC^{\times},
\qquad
q_{\DK} \;:=\; \exp(2 \pi i \hbar)
\;\;\in\;\; \bbC^{\times}.
\]
The two parameters obey the following six identities.
\begin{enumerate}
\item[\textup{(i)}] \emph{Squaring.}
$q_{\KL}^{2} = q_{\DK}$.
\item[\textup{(ii)}] \emph{Logarithmic identity.}
$\log q_{\DK} = 2 \log q_{\KL} = 2 \pi i \hbar$ on the principal
branch of the universal cover $\bbR + i \bbR \to \bbC^{\times}$.
\item[\textup{(iii)}] \emph{r-matrix bridge \textup{(level-prefix
convention: trace-form versus KZ)}.}
The trace-form current-algebra residue
$r_{\mathrm{tr}}(z):=k\,\Omega_{\mathrm{tr}}/z$ and the KZ
connection residue
$r_{\KZ}(z):=\Omega/((k+\hv)z)=\hbar\,\Omega/z$ are different
normalised representatives. The Casimir rescaling
$\Omega=2\hv\,\Omega_{\mathrm{tr}}$ gives
\[
r_{\KZ}(z)=\frac{2\hv}{k+\hv}\,
\frac{\Omega_{\mathrm{tr}}}{z},
\]
so equality with $r_{\mathrm{tr}}(z)$ is not an identity in~$k$.
The trace-form representative carries the level prefix and satisfies
$r_{\mathrm{tr}}(z)|_{k=0}=0$; the KZ representative carries the
KZ coupling and has value $\Omega/(\hv z)$ at $k=0$.
\item[\textup{(iv)}] \emph{R-matrix monodromy.}
For a representation $V_{\lambda}$ of $\fg$, the full pure-braid
monodromy of the Knizhnik--Zamolodchikov connection
$\nabla_{\KZ} = d - \hbar\, \Omega \, dz/z$ on
$\Conf_{2}(\bbC)$ around the diagonal $z_{1} = z_{2}$ equals
$M_{\mathrm{full}} = \exp(2\pi i \hbar \Omega) = q_{\DK}^{\Omega}$;
the half-monodromy generated by the Drinfeld braid generator
$\sigma$ on $\Conf_{2}(\bbC)/S_{2}$ equals
$M_{\mathrm{half}} = \exp(\pi i \hbar \Omega) = q_{\KL}^{\Omega}$;
the universal R-matrix in either convention satisfies
$R_{\KL} = q_{\KL}^{\Omega}$ and $R_{\DK} = q_{\DK}^{\Omega} =
R_{\KL}^{2}$.
\item[\textup{(v)}] \emph{Conformal weights at q-deformation.}
The Sugawara conformal weight $h(V_{\lambda}) =
(\lambda, \lambda + 2\rho)/(2(k + \hv))$ exponentiates
according to $q_{\DK}^{h(V_{\lambda})} = q_{\KL}^{2 h(V_{\lambda})}$
in $\bbC^{\times}$; KL sees half-weights on the cover, DK sees full
weights on the base.
\item[\textup{(vi)}] \emph{Categorical agreement.}
For generic $\hbar$, the braided monoidal category
$\Rep_{q}(\fg)$ is intrinsic to the quantum-group construction
and does not depend on the convention choice. The relabelling
$q \mapsto q^{1/2}$ induces an isomorphism of braided monoidal
categories
$\Rep_{q_{\KL}}(\fg) \xrightarrow{\sim} \Rep_{q_{\DK}^{1/2}}(\fg)$.
At a root of unity, the two conventions change the recorded order
of the parameter and hence the alcove labelling. Simple-object
counts depend on the Lie type, parity, lattice, and chosen
semisimple or small-quantum-group quotient; no universal factor-two
count is asserted here. The invariant statement is the parameter
identity $q_{\KL}^{2}=q_{\DK}$ on the common cover.
\end{enumerate}
\end{theorem}

\begin{proof}
\textbf{(i)} Direct: $q_{\KL}^{2} = \exp(\pi i \hbar)^{2} =
\exp(2 \pi i \hbar) = q_{\DK}$.

\textbf{(ii)} Apply the principal-branch logarithm to (i).

\textbf{(iii)} The trace form $(\cdot, \cdot)_{\mathrm{tr}}$ and
the normalised Killing form $(\cdot, \cdot)$ are related by
$(\cdot, \cdot) = 2\hv (\cdot, \cdot)_{\mathrm{tr}}$, the
standard normalisation fixed by long-root squared length two; the
Casimirs satisfy $\Omega = 2 \hv \, \Omega_{\mathrm{tr}}$.
Substitution gives the displayed expression for $r_{\KZ}$. The two
representatives are used in different calculations:
$r_{\mathrm{tr}}$ records the trace-form current-algebra OPE with
the level as a prefactor, while $r_{\KZ}$ records the KZ connection
residue with coupling $(k+\hv)^{-1}$. Their comparison is therefore
a normalization bridge, not an equality of rational functions in~$k$.

\textbf{(iv)} The KZ connection $\nabla_{\KZ} = d - \hbar
\Omega \, dz/z$ on $\Conf_{2}(\bbC)$ has residue $\hbar \Omega$
at $z = 0$. By the residue calculus, the pure-braid monodromy
along a small loop encircling the diagonal once is $\exp(2\pi i
\hbar \Omega) = q_{\DK}^{\Omega}$. The half-monodromy is the
path interpolating between $(z_{1}, z_{2})$ and the swap
$(z_{2}, z_{1})$; it traverses half the loop and contributes
$\exp(\pi i \hbar \Omega) = q_{\KL}^{\Omega}$. The
Drinfeld--Kohno theorem identifies $M_{\mathrm{full}}$ with the
universal R-matrix of $U_{q_{\DK}}(\fg)$ acting on $V \otimes V$
and $M_{\mathrm{half}}$ with the universal R-matrix of
$U_{q_{\KL}}(\fg)$; their relation $M_{\mathrm{full}} =
M_{\mathrm{half}}^{2}$ is the squaring identity \textup{(i)}
lifted to operators.

\textbf{(v)} The conformal weight $h(V_{\lambda})$ arises from the
Sugawara $L_{0}$ eigenvalue. The phase rotation
$M_{\mathrm{full}} \cdot |V_{\lambda}\rangle = \exp(2\pi i
h(V_{\lambda})) |V_{\lambda}\rangle$ agrees with
$q_{\DK}^{h(V_{\lambda})} |V_{\lambda}\rangle$ in the DK
convention; rewriting via $q_{\DK} = q_{\KL}^{2}$ gives
$q_{\KL}^{2 h(V_{\lambda})}$.

\textbf{(vi)} The braided monoidal category $\Rep_{q}(\fg)$ at
generic $q$ is determined by the universal R-matrix and its
intertwiner data, both of which are intrinsic to the
quantum-group construction. The relabelling $q \mapsto q^{1/2}$
identifies the generic braided categories after passing to the
common square-root cover. At roots of unity, the recorded order of
the parameter and the alcove labels change. The object count then
depends on the chosen root datum and quotient category, so the
theorem records only the parameter relation.
\end{proof}

\begin{remark}[Chern--Simons interpretation]
\label{rem:q-bridge-cs}
The KL convention is natural for Chern--Simons / Wess--Zumino--Witten
because Witten 1989 writes the action as $S_{\mathrm{CS}} =
(k/(4\pi)) \int \mathrm{Tr}(A \wedge dA + (2/3) A \wedge A \wedge
A)$ and obtains $q = \exp(\pi i / (k + \hv)) = q_{\KL}$. The factor
of $\pi$, not $2\pi$, is the half-monodromy of a Wilson line
around itself; the full monodromy of one Wilson line around
another is $q_{\KL}^{2} = q_{\DK}$. The KL parameter is the
self-monodromy parameter; the DK parameter is the
mutual-monodromy parameter.
\end{remark}

\begin{remark}[Kazhdan--Lusztig periodicity]
\label{rem:q-bridge-periodicity}
At admissible level $k = -\hv + p/q$ with $\gcd(p, q) = 1$ and
$p \geq \hv$ for simply-laced $\fg$, $q_{\KL} = \exp(\pi i q/p)$
is a $2p$-th root of unity, while $q_{\DK} = \exp(2\pi i q/p)$
is a $p$-th root of unity. The Kazhdan--Lusztig equivalence is
stated at $q_{\KL}$ and the Drinfeld--Kohno theorem is stated at
$q_{\DK}$. The squaring identity changes the recorded root order
from $2p$ to $p$. It does not by itself determine the number of
simple objects in a root-of-unity quotient; that number depends on
the root datum, parity, lattice convention, and the quotient category.
\end{remark}

\section{Internal consistency: five derived quantities}
\label{sec:q-bridge-consistency}

We verify that the bridge is closed under all derived
constructions used in the manuscript.

\paragraph{R-matrix.}
Both conventions encode the same universal R-matrix as an element
of $(U_{q}(\fg) \otimes U_{q}(\fg))^{\mathrm{completed}}$. In KL,
$R_{\KL} = q_{\KL}^{\Omega} \cdot (\mathrm{correction})$; in DK,
$R_{\DK} = q_{\DK}^{\Omega} \cdot (\mathrm{correction})$. The
correction terms are q-series in the Lie algebra side, independent
of convention, built from $q_{\KL} = q_{\DK}^{1/2}$ on the
diagonal Cartan. The squaring relation is a parameter relation on
the common cover; the full universal R-matrix is compared by
rewriting its q-series in that parameter, not by literally squaring
the complete R-matrix.

\paragraph{KZ monodromy.}
The pure-braid full monodromy is $\exp(2\pi i \hbar \Omega) =
q_{\DK}^{\Omega}$ by Cauchy residues. The half-monodromy is
$q_{\KL}^{\Omega}$. No ambiguity.

\paragraph{Conformal weights.}
The conformal weight $h(V_{\lambda})$ is convention-independent
(computed from Sugawara $L_{0}$ with denominator $k + \hv$). The
phase $q^{2h}$ exponentiates to $\exp(2\pi i \cdot 2 h \cdot
\hbar)$ which coincides with $q_{\DK}^{h}$ and with
$q_{\KL}^{2h}$ in $\bbC^{\times}$.

\paragraph{Verlinde fusion.}
The Verlinde formula expresses fusion coefficients $N^{k}_{ij}$
through the modular S-matrix entries $S_{ij} \propto \exp(\pm\pi
i (\lambda + \rho) \cdot (\mu + \rho)/(k + \hv)) =
q_{\KL}^{\pm(\lambda + \rho) \cdot (\mu + \rho)}$. Recasting in
DK gives $q_{\DK}^{\pm(\lambda + \rho) \cdot (\mu + \rho)/2}$.
Both produce the same fusion coefficients.

\paragraph{Modular T-matrix.}
The modular T-matrix is $T = q_{\DK}^{h - c/24}$ in the DK
convention and $T = q_{\KL}^{2(h - c/24)}$ in the KL convention.
The modular invariants $T^{N} S^{M}$ match.

The conclusion of \S\ref{sec:q-bridge-consistency} is that all
five derived quantities are consistent under the bridge: there is
no derived object in the manuscript for which the two conventions
disagree. The bridge is closed under every operation Vol~I
performs.

\section{The bridge as KZ cocycle}
\label{sec:q-bridge-cocycle}

The structural content of \eqref{eq:q-bridge} is the index-2
inclusion of the pure-braid subgroup into the full Artin braid
group on two strands.

\begin{theorem}[Q-bridge as Z/2-cover cocycle]
\label{thm:q-bridge-cocycle}
\ClaimStatusProvedHere
Let $M = M_{\KZ}(\fg, k)$ denote the universal monodromy local
system on $\Conf_{2}(\bbC)/S_{2}$ of the KZ connection
$\nabla_{\KZ} = d - \hbar\, \Omega \, dz/z$ at non-critical level
$k$, and let $B_{2} = \pi_{1}(\Conf_{2}(\bbC)/S_{2})$ denote the
Artin braid group on two strands. There is a canonical short
exact sequence
\[
1 \;\to\; \langle \sigma^{2} \rangle \;\to\; B_{2} \;\to\; S_{2}
\;\to\; 1,
\]
where $\sigma$ is the Drinfeld braid generator and
$\sigma^{2}$ generates the pure-braid subgroup $P_{2} =
\pi_{1}(\Conf_{2}(\bbC))$. The KZ monodromy
$\rho_{\KZ} \colon B_{2} \to \mathrm{Aut}(V \otimes V)$ factors
through this extension via $\rho_{\KZ}(\sigma^{2}) = q_{\DK}^{\Omega}$
(image of $P_{2}$, full pure-braid monodromy) and
$\rho_{\KZ}(\sigma) = q_{\KL}^{\Omega}$ (image of the swap class,
half-monodromy). The Drinfeld--Kohno theorem identifies
$\rho_{\KZ}$ with the universal R-matrix representation of
$U_{q}(\fg)$; the convention choice is the choice of which
level of the extension to take as the parameter base. The KL
convention takes $B_{2}$ as the base (the cover); the DK
convention takes $\langle \sigma^{2} \rangle = P_{2}$ as the base
(the quotient by the order-two framing). The squaring identity
$q_{\KL}^{2} = q_{\DK}$ realises the index-2 inclusion
$\langle \sigma^{2} \rangle \subset B_{2}$ algebraically.
\end{theorem}

\begin{proof}
The double cover $\Conf_{2}(\bbC) \to \Conf_{2}(\bbC)/S_{2}$ is
the universal $\bbZ/2$-cover, with deck group generated by the
swap $\sigma \in B_{2}/P_{2}$. The KZ connection lifts to the
cover; the lifted holonomy along $\sigma$ is $\exp(\pi i \hbar
\Omega) = q_{\KL}^{\Omega}$. The original holonomy on the base
along $\sigma^{2}$ is $\exp(2\pi i \hbar \Omega) =
q_{\DK}^{\Omega}$. The convention choice is the choice of working
on the cover (KL, fine framed datum) or the base (DK, coarse
unframed datum); the cover tracks the framing of a single strand
around itself, corresponding to the self-monodromy (topological
spin / framing anomaly) of a Wilson line in Chern--Simons. The
DK convention forgets this anomaly; the KL convention retains
it.
\end{proof}

\begin{remark}[The square root is the lift]
\label{rem:q-bridge-square-root}
The square root in $q_{\KL} = q_{\DK}^{1/2}$ is the algebraic
shadow of the topological double cover. The Drinfeld associator
$\Phi_{\KZ}(t_{12}, t_{23}) \in U(\mathfrak{t}_{3})$ governing
the KZ connection on $\Conf_{3}(\bbC)$ is written in
the DK convention because Drinfeld's 1989 paper uses $2\pi i$;
the KL realisation of the same associator is a square root,
choosing a branch. The cocycle condition for the KZ connection
on $\Conf_{n}(\bbC)$, which makes the monodromy a representation
of the braid group rather than just a cochain, is the squaring
relation \eqref{eq:q-bridge}. The convention choice is a gauge
choice: KL works on the universal cover; DK works on the base.
The bridge \eqref{eq:q-bridge} is the gauge transformation
between them.
\end{remark}

\section{Scope of the convention bridge}
\label{sec:q-bridge-propagation}

\(\eqref{eq:q-bridge}\) is an identity between two parameter
normalizations.  It does not identify the trace-form current residue with
the KZ residue, nor does it by itself identify full universal
\(R\)-matrices or root-of-unity quotient categories.  Those comparisons
require the additional hypotheses stated in
Theorem~\ref{thm:q-convention-bridge-main}.

\subsection{Formal substitution}
\label{sec:q-bridge-protocol}

For an expression already defined as a function of the quantum parameter,
the two notations are related on a chosen square-root cover by
\[
 q_{\DK}=q_{\KL}^{2},\qquad
 q_{\KL}=q_{\DK}^{1/2}
\]
as a formal change of variable.  No further comparison theorem is asserted
by this substitution.

\section{Summary}
\label{sec:q-bridge-summary}

The proved parameter identity is \(q_{\KL}^{2}=q_{\DK}\).  The surrounding
operator, residue, conformal-weight, and categorical assertions retain the
hypotheses and scope of Theorem~\ref{thm:q-convention-bridge-main}.  In
particular, the square-root notation alone determines neither an alcove nor
a root-of-unity quotient category.
